\par\noindent\begin{minipage}{\linewidth}
% !TeX root = ../教材题库_全部条目_转写工作稿.tex
% Source: XJ-B1-C04-TH-P106-Y0720
% 第4章 / 4.2 指数函数 / 教材 p.106 / PDF 第113页
% 已对照原页转写题干及公式；待教师复核；未制作答案
% 原题及相邻原解截图（核对用）：../题目原图/XJ-B1-C04-TH-P106-Y0720-01.webp
% 本文件只在转写完成后作为题干使用；OCR 条目尚不是可靠题干。
\begin{bankitem}{p.106 思考}{第4章 / 4.2 指数函数 / 教材 p.106 / PDF 第113页}{XJ-B1-C04-TH-P106-Y0720}{已对照原页转写题干及公式；待教师复核；未制作答案}
\subq{1}将一张纸对折一次得到2层，对折两次得到4层。对折 $x$ 次后，纸的层数 $y$ 与 $x$ 之间有什么关系？
\subq{2}《庄子》中说：“一尺之棰，日取其半，万世不竭。”若每天截取木棰的一半，剩余木棰的长度 $y$（尺）与截取的天数 $x$ 之间有什么关系？
\subq{3}上述两个关系有什么共同特征？
\end{bankitem}
\end{minipage}\par\medskip
